Просвітлення, про яке я кажу, — це не просто здійснення, не просто відкриття власної справжньої природи. Це лише початок — точка входу у внутрішню революцію. Здійснення не гарантує цієї революції; воно лише робить її можливою.

\section*{Внутрішня революція}

Що це за внутрішня революція? Передусім революція не статична; вона жива, триває й не припиняється. Її не можна схопити чи втиснути в якусь понятійну модель. І немає жодного шляху до цієї внутрішньої революції, бо вона непередбачувана й некерована і має власне життя. Ця революція — це відрив від старих, повторюваних, мертвих структур думки й сприйняття, у яких людство опинилося в пастці. Здійснення останньої Дійсності — це безпосереднє й раптове екзистенційне пробудження до власної справжньої природи, яке відчиняє двері до можливості внутрішньої революції. Така революція вимагає безперервного спорожнення старих структур свідомості й народження живої, плинної інтелігентності. Ця інтелігентність перебудовує всю вашу істоту — тіло, розум і сприйняття. Вона вивільняє розум зі старих структур, укорінених у сукупності людської свідомості. Якщо не звільнитися від старих зумовлених структур людської свідомості, людина й далі перебуває у в'язниці.

Пробудження до власної справжньої природи не обов'язково означає, що відбуватиметься безперервна революція в тому, як людина сприймає життя, діє й відповідає на нього. Мить пробудження показує нам, що остаточно істинне й реальне, а також відкриває глибшу можливість того, як жити з нерозділеного й незумовленого стану буття. Але мить пробудження не гарантує цієї глибшої можливості — це можуть засвідчити багато з тих, хто пережив духовне пробудження. Пробудження відчиняє внутрішні двері до глибокої внутрішньої революції, але аж ніяк не гарантує, що вона відбудеться. Відбудеться вона чи ні — залежить від багатьох чинників, але жоден не є важливішим і життєвішим за щирий і недвозначний намір Істини понад усе. Від цього щирого наміру до Істини зрештою залежить усе духовне зростання, особливо коли він перевершує всі особисті вподобання, програми й цілі.

Ця внутрішня революція — це пробудження інтелігентності, народженої не з розуму, а з внутрішньої тиші розуму, яка одна здатна викорінити всі старі структури свідомості людини. Якщо ці структури не викорінити, не буде ні творчої думки, ні творчої дії, ні творчої відповіді. Без внутрішньої революції ніщо нове й свіже не зможе розквітнути. За відсутності цієї революції розквітне лише старе, повторюване, зумовлене. Але наш потенціал лежить за межами відомого, за межами структур минулого, за межами всього, що встановило людство. Наш потенціал може розквітнути лише тоді, коли ми більше не впіймані у впливі й обмеженнях відомого. За межами царини розуму, за межами обмежень зумовленої свідомості людства лежить те, що можна назвати священним. І саме зі священного народжується нова, плинна свідомість, яка стирає старе й оживляє розквіт живого й нерозділеного вияву буття. Такий вияв не особистий і не безособовий, не духовний і не мирський, а радше течія й розквіт існування за межами всіх уявлень про себе.
